\section{Modelo y ley de control propuesta}

\subsection{Estado del sistema}

El estado minimo contiene robots, cargas, obstaculos y grafo:
\[
  x = \{(q_i,\theta_i,b_i,c_i)_{i=1}^{N},
  (q_k^L,w_k^{\rm dem},D_k,\ell_k,a_k)_{k=1}^{M}, G(t)\}.
\]
Cada robot observa informacion local: vecinos, cargas descubiertas, estimacion de
ocupacion, distancia, bateria y margen de seguridad. Para cargas no puntuales, $\ell_k$
y $a_k$ representan largo y ancho de la huella rectangular usada por la formacion rigida.

\subsection{Lectura escalonada del modelo}

La formulacion se construye de menor a mayor exigencia fisica. En M1, el robot solo
elige carga:
\[
  y_i \in \Delta_M,\qquad
  C_k(t)=\{i:\argmax_{\ell} y_{i\ell}(t)=k\}.
\]
En el caso homogeneo M1.1, una carga queda servida si la coalicion alcanza cardinalidad
minima:
\[
  |C_k| \ge n_k .
\]
En M1.2--M1.4, la cardinalidad se reemplaza por capacidad efectiva:
\[
  S_k(C_k,t)=\sum_{i\in C_k} c_{ik}(t),\qquad
  c_{ik}(t)=c_i\,\eta_{ik}(q_i,b_i,G(t)),
  \qquad S_k(C_k,t)\ge D_k .
\]
El factor $\eta_{ik}$ resume la perdida por distancia, orientacion, bateria, saturacion,
contacto inutil o desconexion. En M3, el robot no solo elige carga sino tambien rol o
slot:
\[
  z_{ikh}\in\{0,1\},\qquad
  \sum_{k,h} z_{ikh}\le 1,\qquad
  C_{kh}=\{i:z_{ikh}=1\}.
\]
La implementacion actual usa esta idea de manera parcial: Smith-QR asigna la carga y la
formacion deriva el slot por ranking dentro de la coalicion. La optimizacion completa
$z_{ikh}$ queda como extension natural.

\subsection{Comunicacion y memoria local}

La comunicacion se modela mediante un grafo $G(t)$ dependiente de rango, perdida y
descubrimiento local. La condicion nominal para el analisis es conectividad a trozos,
pero los experimentos fuerzan regimenes degradados como R3 para medir que ocurre cuando
Smith pierde informacion global. Smith-QR no asume broadcast perfecto: cada robot
mantiene memoria de descriptores de carga, estima ocupacion por consenso dinamico y usa
compromiso temporal para evitar conmutaciones espurias.

\subsection{Smith-QR}

Smith-QR conserva la dinamica poblacional de Smith, pero agrega los mecanismos necesarios
para que el resultado sea ejecutable:
\begin{itemize}[leftmargin=2em]
  \item memoria de descriptores de carga bajo comunicacion local;
  \item estimacion por consenso dinamico de ocupacion;
  \item clearing entero para cerrar quorums utiles;
  \item compromiso en ruta para reducir chattering;
  \item guardias de quorum para no declarar transporte sin coalicion fisica.
\end{itemize}

\subsection{Campo unico continuo}

La ley deseada para el robot $i$ se expresa como suma suave de campos:
\[
  u_i =
  -\nabla_{q_i}\Phi_i^{\rm tarea}
  -\nabla_{q_i}\Phi_i^{\rm formacion}
  -\nabla_{q_i}\Phi_i^{\rm obstaculo}
  -\nabla_{q_i}\Phi_i^{\rm bateria}
  + u_i^{\rm wrench}.
\]
El comportamiento emergente depende del termino dominante, no de un arbol de \texttt{if}
centrado en modos. Si la bateria cae, el campo de retornabilidad gana peso; si una carga
requiere torque, el termino wrench favorece contactos utiles; si el quorum se cierra, el
deficit baja y el reclutamiento se relaja.

\subsection{Formacion rigida de transporte}

La asignacion Smith-QR produce una coalicion $C_k$, pero no basta con saber cuantos robots
hay. Para transportar una carga extendida se define una formacion rigida solidaria al
marco de la carga:
\[
  s_j = (r_j,n_j), \qquad r_j \in \RR^2,\quad \norm{n_j}=1,
\]
con $r_j$ como punto de contacto relativo al centro de masa y $n_j$ como direccion de
fuerza admisible. En la implementacion V7, los slots se generan a partir de la huella
rectangular de la carga $(\ell_k,a_k)$: esquinas y lados permiten producir fuerza
longitudinal, fuerza lateral y torque planar. El robot asignado con rango $j$ en la
coalicion apunta a $q_k^L+r_j$ y usa $n_j$ como direccion mecanica preferente.

Esta decision convierte la formacion en parte del modelo de carga, no en un circulo
visual arbitrario alrededor del objetivo. Tambien permite representar cargas largas,
desbalanceadas o con restricciones de contacto sin cambiar la interfaz del asignador.

\subsection{Juego vectorial de wrench}

Para una carga plana, cada contacto aporta fuerza y torque:
\[
  w_j =
  \begin{bmatrix}
  n_{jx}\\ n_{jy}\\ r_{jx}n_{jy}-r_{jy}n_{jx}
  \end{bmatrix}\lambda_j,
  \qquad 0 \le \lambda_j \le \bar f_j,
\]
donde $\lambda_j$ es el esfuerzo escalar elegido por el robot del slot $j$ y $\bar f_j$
es su limite individual. El juego vectorial se formula como una mejor respuesta
cooperativa en espacio wrench:
\[
  \lambda_C^\star =
  \argmin_{0 \le \lambda \le \bar f}
  \norm{G_C\lambda - w_k^{\rm dem}}_2^2.
\]
La utilidad marginal de cada robot se interpreta como su reduccion del residual comun.
Asi, una coalicion es fisicamente aceptable si el residual
$\norm{G_C\lambda_C^\star-w_k^{\rm dem}}$ y la saturacion por robot quedan bajo umbrales
definidos. Esta es la diferencia entre ``hay suficientes robots'' y ``la carga puede
moverse con fuerza y torque viables''.
